% Báo cáo cuối của đề tài dlinear-objective-lab (NHIEM_VU_8, Phần B).
% Preamble chép nguyên từ report/template/template.tex (template AI VIET NAM, AIO2026), phần trước \begin{document};
% mọi chỗ đổi có chú thích "NV8". Biên dịch: xem report/README.md.
% documentclass nằm trong main.tex (NV8: để Overleaf không chọn nhầm file này làm file chính)
\usepackage{makecell}
\usepackage{hyphenat}
\usepackage[utf8]{vietnam}
\usepackage[vietnamese=nohyphenation]{hyphsubst}
\usepackage[vietnamese]{babel}
\usepackage[table,dvipsnames]{xcolor} % Required for row coloring
\usepackage{booktabs} % Required for nicer horizontal lines
\usepackage{ragged2e}
\usepackage{caption}
\usepackage{subcaption}
\usepackage{float}
\usepackage{amsmath,amssymb,amsfonts}
\usepackage{graphicx}
\usepackage{adjustbox}
\usepackage{multicol}
\usepackage[inline,shortlabels]{enumitem}
\usepackage{setspace}
\usepackage{longtable}
\usepackage{bm} 
\usepackage[inkscapelatex=false]{svg}
\usepackage{titling}
\usepackage{tabularx}
\usepackage[T1]{fontenc}
\usepackage{lmodern,mathrsfs}
% \usepackage[draft]{hyperref}
\usepackage{hyperref}
\usepackage{xparse}
\usepackage{fontawesome5}
\usepackage[most]{tcolorbox}
\usepackage{tikz}
\usepackage[framemethod=TikZ]{mdframed}
\usetikzlibrary{matrix, fit, shadows, positioning, arrows.meta, calc, shapes.geometric, backgrounds, shapes.symbols}
\usepackage{paracol}
\usepackage{bbm}
\usepackage{pgfplots}
\usepackage{siunitx}
\usepackage{multirow}
%  biblatex (thay cho IEEEtran.bst) ------------------------------------
\usepackage[
    backend=bibtex,
    style=ieee,
    citestyle=numeric-comp,
    maxnames=1,     % chỉ hiển thị 1 tên tác giả, phần còn lại thành ``et al."
    minnames=1,
    uniquelist=false,
    hyperref=false,
    language=english  % NV8: chuỗi biblatex tiếng Anh, thay cho \selectlanguage{english} bên dưới
]{biblatex}
\usepackage{forest}
\addbibresource{refs.bib}  % NV8: đổi tên file .bib
\DefineBibliographyStrings{english}{
  andothers = {et\addabbrvspace al\adddot}
}
% NV8: bỏ \AtBeginBibliography{\selectlanguage{english}}: nó đổi font sang T1, chân trang có chữ Việt ("Hồng Phúc") làm lỗi \ocircumflex unavailable in encoding T1


% -----------------------------------------------------------------
\AtBeginDocument{\renewcommand\refname{}}

\usepackage[toc,page]{appendix}
\renewcommand{\appendixtocname}{Phụ lục}      
\renewcommand{\appendixpagename}{Phụ lục}

\usepackage{xurl}
\setlength{\headheight}{14.5pt}
\setlength{\topmargin}{-.5in} \setlength{\textheight}{9.25in}
\setlength{\oddsidemargin}{-0.1in} \setlength{\textwidth}{6.8in}

\usepackage{tocloft}
\renewcommand{\cfttoctitlefont}{\Huge\bfseries}
\setlength{\cftaftertoctitleskip}{2em}
\setlength{\cftbeforesubsecskip}{7pt} % giãn cách subitem mục lục
\renewcommand{\cftsecleader}{\cftdotfill{\cftdotsep}}
\setlength{\cftsecnumwidth}{3em}
\renewcommand{\cftsecaftersnum}{\quad}
\renewcommand{\cftsubsecleader}{\cftdotfill{\cftdotsep}}
\setlength{\cftsubsecnumwidth}{4em}
\renewcommand{\cftsubsecaftersnum}{\quad}
\renewcommand{\cftsubsubsecleader}{\cftdotfill{\cftdotsep}}
\setlength   {\cftsubsubsecnumwidth}{5em}  
\renewcommand{\cftsubsubsecaftersnum}{\quad}

\PassOptionsToPackage{vietnamese.licr}{babel}
\usepackage{vipythonhighlight}
\usepackage{tvietlistings}
\newcounter{mycounter}
\newcommand\showmycounter{\stepcounter{mycounter}\themycounter}
\usepackage{lipsum}
\newcommand\showlips{\stepcounter{mycounter}\lipsum[\value{mycounter}]}

%--------------------------------------------------------------------
\usepackage{framed}
\usepackage{fancyhdr}
\usepackage{upquote}
\usepackage{listings}	 
	    

\usepackage[toc, nonumberlist, style=super]{glossaries}
\makenoidxglossaries

\newglossarystyle{clean3col}{%
    % 1) Table layout: 3 columns
    \renewenvironment{theglossary}%
    {%
      \begin{longtable}{@{}
        >{\setlength{\parindent}{0pt}\RaggedRight\arraybackslash}p{0.2\linewidth}
        >{\setlength{\parindent}{0pt}\RaggedRight\arraybackslash}p{0.23\linewidth}
        >{\setlength{\parindent}{0pt}\RaggedRight\arraybackslash}p{0.52\linewidth}
      @{}}%
    }%
    {\end{longtable}}%

    % 2) Header row
    \renewcommand{\glossaryheader}{%
        \bfseries Thuật ngữ &
        \bfseries Tạm dịch &
        \bfseries Mô tả \\
        \midrule
        \endhead
    }%

    % 3) Each entry row
    \renewcommand{\glossentry}[2]{%
        \itshape \glstarget{##1}{\glossentryname{##1}} &
        \glsentryuseri{##1} &
        \glossentrydesc{##1} \\
        \addlinespace[6pt]
    }%
}


% Color --------------------------------------------------------------------
\usepackage{epigraph}
\usepackage{titlesec}

\definecolor{codebackground}{RGB}{250, 250, 250}
\definecolor{codecomment}{RGB}{34, 139, 34}
\definecolor{codekeyword}{RGB}{0, 0, 255}
\definecolor{codestring}{RGB}{255, 0, 0}
\definecolor{codeidentifier}{RGB}{0, 0, 0}
\definecolor{codelinenumbers}{RGB}{90, 90, 90}

\definecolor{codegreen}{rgb}{0,0.6,0}
\definecolor{codegray}{rgb}{0.5,0.5,0.5}
\definecolor{codepurple}{rgb}{0.58,0,0.82}
\definecolor{backcolour}{rgb}{0.95,0.95,0.92}

\definecolor{aivnGreen}{RGB}{52,79,30}  
\definecolor{aivnCodeOutputColor}{RGB}{11,83,148}  
\definecolor{captiongray}{gray}{0.15} 
\definecolor{captionbg}{RGB}{250,250,250}

\definecolor{foldercolor}{RGB}{0, 102, 204}
\definecolor{filecolor}{RGB}{104, 104, 104}
\definecolor{scriptcolor}{RGB}{180, 50, 50}
\definecolor{pythoncolor}{RGB}{0, 128, 128}

% --- Extra Brand colors for the table ---
\definecolor{brandIndigo}{HTML}{333D91} 
\definecolor{rowLavender}{HTML}{EEECFA}
\setlength{\parskip}{0.3em}

% Table
\newtcolorbox{captionbelowboxaivn}{
    colback=aivnGreen,
    colframe=aivnGreen,
    boxrule=0pt,
    arc=4pt,
    left=4pt,
    right=4pt,
    top=2pt,
    bottom=2pt,
    width=\linewidth,
    boxsep=4pt,
    halign=center,
    % fontupper=\bfseries\color{white}
    fontupper=\color{white}
}

% Figure
\newtcolorbox{figurecardcaptionaivn}{
  colback=aivnGreen,        % Màu nền (background color)
  colframe=aivnGreen,       % Màu viền box (border color)
  boxrule=0pt,              % Độ dày viền box (0pt = không viền)
  arc=2pt,                  % Độ bo góc (2pt là bo nhẹ, tăng để bo tròn hơn)
  top=3pt,                  % Padding phía trên nội dung
  bottom=3pt,               % Padding phía dưới nội dung
  left=10pt,                % Padding bên trái
  right=10pt,               % Padding bên phải
  width=\linewidth,         % Chiều rộng bằng toàn bộ dòng
  fontupper=\small\fontsize{11pt}{13pt}\selectfont\color{white} 
}

% Tạo counter riêng
\newcounter{hinh}
% Định nghĩa macro caption Hình
\newcommand{\capfig}[1]{%
  \refstepcounter{hinh}%
  % \textbf{Hình~\thehinh:}~#1
    Hình~\thehinh:~#1
}

\newtcolorbox{figurecaptionmodern}{
  colback=captionbg,
  colframe=aivnGreen,
  boxrule=1.5pt,
  bottomrule=2pt,  
  toptitle=0pt,
  bottomtitle=0pt,
  left=6pt,
  right=6pt,
  top=4pt,
  bottom=2pt,
  width=\linewidth,
  fontupper=\small\fontsize{11pt}{13pt}\selectfont\color{black},
  halign=center
}

% Code Block -----------------------------------------------------------------
\newtcolorbox{aivncodebox}[1][]{%
  enhanced,
  colback=gray!10,
  colframe=aivnGreen,
  coltitle=white,
  colbacktitle=aivnGreen,
  title=#1,
  fonttitle=\bfseries\sffamily,
  listing only,
  breakable,
  listing engine=listings,
  listing options={
    language=Python,
    basicstyle=\ttfamily\footnotesize,
    commentstyle=\color{ForestGreen},
    keywordstyle=\color{blue},
    stringstyle=\color{red!70!black},
    numbers=left, upquote=true,
    numberstyle=\tiny\color{gray},
    numbersep=5pt,
    breaklines=true,
    showstringspaces=false,
    xleftmargin=1.5em,
    framexleftmargin=1.5em
  }
}

% Code Output Block -----------------------------------------------------------------
\newtcolorbox{aivnCodeOutput}[1][]{%
  enhanced,
  colback=gray!10,
  colframe=aivnCodeOutputColor,
  coltitle=white,
  colbacktitle=aivnCodeOutputColor,
  title=#1,
  fonttitle=\bfseries\sffamily,
  listing only,
  breakable,
  listing engine=listings,
  listing options={
    language=Python,
    basicstyle=\ttfamily\footnotesize,
    commentstyle=\color{ForestGreen},
    keywordstyle=\color{blue},
    stringstyle=\color{red!70!black},
    numbers=left,
    numberstyle=\tiny\color{gray},
    numbersep=5pt,
    breaklines=true,
    showstringspaces=false,
    xleftmargin=1.5em,
    framexleftmargin=1.5em
  }
}

% Last argument allows more tcolorbox options to be added

% --- Guideline/Tip box ---
\newtcolorbox{guidelinebox}[2][]{
    enhanced, colback=yellow!10!white, colframe=orange!90!black, boxrule=1pt, arc=4pt,
    fonttitle=\bfseries, attach boxed title to top left={yshift=-2mm, xshift=3mm},
    colbacktitle=orange!90!black, coltitle=white,
    title=\faLightbulb\hspace{1mm} #2, #1
}

\newtcolorbox{notebox}[1][]{
  enhanced,  colback=yellow!10, colframe=orange!80, boxrule=1pt, 
  arc=4pt, fontupper=\small, attach boxed title to top left={yshift=-2mm, xshift=3mm},
  colbacktitle=orange!80,  coltitle=white, fonttitle=\bfseries,  
  title=\faInfoCircle~\textbf{Lưu ý},  #1
}

\newtcolorbox{infobox}[1]{%
    enhanced, 
    breakable, 
    colback=green!3, 
    colframe=green!3, 
    arc=5pt,
    borderline west={3pt}{0pt}{aivnGreen},
    attach boxed title to top left={ 
        yshift*=-\tcboxedtitleheight/2, 
        xshift=-0.1em 
    },
    colbacktitle=aivnGreen, %
    coltitle=white, 
    fonttitle=\bfseries, 
    boxed title style={ 
        arc=4pt, 
        top=2pt, 
        bottom=2pt,
        left=8pt,
        right=8pt,
    },
    title={\faInfoCircle\quad #1}
}
\newtcolorbox{examplebox}[1]{%
    enhanced,
    breakable,
    colback=brandIndigo!5,      
    colframe=brandIndigo!5,     
    arc=5pt,
    borderline west={3pt}{0pt}{brandIndigo}, 
    attach boxed title to top left={
        yshift*=-\tcboxedtitleheight/2,
        xshift=-0.1em
    },
    colbacktitle=brandIndigo,   
    coltitle=white,
    fonttitle=\bfseries,
    boxed title style={
        arc=4pt,
        top=2pt,
        bottom=2pt,
        left=8pt,
        right=8pt,
    },
    title={\faLightbulb\quad #1} 
}
% Prompt Block ---------------------------------------------------------------
\definecolor{promptBlue}{HTML}{2F5D8C}
\definecolor{promptBg}{HTML}{F6F9FC}
\definecolor{promptBorder}{HTML}{B7C9DC}

% Dùng cho prompt đã viết theo cú pháp LaTeX bình thường
\newtcolorbox{promptbox}[1][Prompt]{%
    enhanced,
    breakable,
    colback=promptBg,
    colframe=promptBorder,
    coltitle=white,
    colbacktitle=promptBlue,
    fonttitle=\bfseries\sffamily,
    title={#1},
    boxrule=0.8pt,
    arc=4pt,
    left=8pt,
    right=8pt,
    top=2pt,
    bottom=1pt,
    before skip=8pt,
    after skip=8pt,
    borderline west={3pt}{0pt}{promptBlue},
    fontupper=\small\sffamily,
}

% Dùng cho prompt raw: giữ nguyên xuống dòng, không highlight keyword
\newtcblisting{promptboxraw}[1][Prompt]{%
    enhanced,
    breakable,
    listing only,
    listing engine=listings,
    colback=promptBg,
    colframe=promptBorder,
    coltitle=white,
    colbacktitle=promptBlue,
    fonttitle=\bfseries\sffamily,
    title={#1},
    boxrule=0.8pt,
    arc=4pt,
    left=8pt,
    right=8pt,
    top=-4pt,
    bottom=-4pt,
    before skip=4pt,
    after skip=4pt,
    borderline west={3pt}{0pt}{promptBlue},
    listing options={
        language={},
        basicstyle=\small\ttfamily,
        columns=fullflexible,
        keepspaces=true,
        breaklines=true,
        showstringspaces=false,
        numbers=none,
        frame=none,
        keywordstyle=\color{black},
        commentstyle=\color{black},
        stringstyle=\color{black},
        identifierstyle=\color{black}
    }
}


%--------------------------------------------------------------------
\lstdefinestyle{mystyle}{
    backgroundcolor=\color{codebackground},
    commentstyle=\color{codecomment},
    keywordstyle=\color{codekeyword},
    stringstyle=\color{codestring},
    identifierstyle=\color{codeidentifier},
    basicstyle=\ttfamily\footnotesize,
    breakatwhitespace=false,
    breaklines=true,
    captionpos=b,
    keepspaces=true,
    numbers=left,
    numberstyle=\tiny\color{codelinenumbers}, 
    numbersep=5pt,
    showspaces=false,
    showstringspaces=false,
    showtabs=false,
    tabsize=2
}

\lstset{style=mystyle}

\hypersetup{
    colorlinks=true,
    linkcolor=black,
    filecolor=magenta,    
    citecolor=green,
    urlcolor=blue,
    pdftitle={dlinear-objective-lab},  % NV8: đổi tiêu đề PDF
    pdfpagemode=FullScreen,
}

%--------------------------------------------------------------------
\newcounter{choice}
\renewcommand\thechoice{($\alph{choice}$)}
\newcommand\choicelabel{\thechoice}

%--------------------------------------------------------------------
\newenvironment{choices}%
  {\list{\choicelabel}%
     {\usecounter{choice}\def\makelabel##1{\hss\llap{##1}}%
       \settowidth{\leftmargin}{W.\hskip\labelsep\hskip 0 em}%
       \def\choice{%
         \item
       } % choice
       \labelwidth\leftmargin\advance\labelwidth-\labelsep
       \topsep=0pt
       \partopsep=0pt
     }%
  }%
  {\endlist}
  
%--------------------------------------------------------------------
\def\CChoice{%
      \choice
      \addanswer{\theenumi}{\thechoice}%
    }
    \let\CChoice\CChoice
%    \par % Uncomment this to have choices always start a new line
   % \let\par\@empty
    % If we're continuing the textbf containing the question,
    % then leave a bit of space before the first choice:
    \ifvmode\else\enskip\fi
    \ignorespaces
  %
  {}
\makeatother
\newbox\allanswers
\setbox\allanswers=\hbox{}
\newcommand{\addanswer}[2]{%
  \global\setbox\allanswers=\hbox{\unhbox\allanswers \quad #1.~#2}%
}
\newcommand{\showanswers}{%
  \vfill
  \begin{center}
    Answers
  \end{center}
  \unhbox\allanswers
}

\makeatletter
\newenvironment{multichoices}[1][1]{%
\begin{multicols}{#1}}{%
\end{multicols}}
\makeatother

\pagestyle{fancy}

\DeclareMathOperator*{\argmax}{arg\,max}
\DeclareMathOperator*{\argmin}{arg\,min}
\DeclareMathOperator{\E}{\mathbb{E}}

% Outline --------------------------------------------------------------------
% Section numbering ----------------------------------------------------------
\renewcommand{\thesection}{\Roman{section}.} 
\renewcommand{\thesubsection}{\thesection\arabic{subsection}.}
\renewcommand{\thesubsubsection}{\thesubsection\arabic{subsubsection}.}

% Paper-like section title sizes --------------------------------------------
% Title paper: 17pt
% Section should be clearly smaller than title but still prominent.
\titleformat{\section}
  {\normalfont\fontsize{16pt}{20pt}\selectfont\bfseries}
  {\thesection}
  {0.85em}
  {}

\titleformat{\subsection}
  {\normalfont\fontsize{14pt}{16pt}\selectfont\bfseries}
  {\thesubsection}
  {0.75em}
  {}

\titleformat{\subsubsection}
  {\normalfont\fontsize{12pt}{15pt}\selectfont\bfseries}
  {\thesubsubsection}
  {0.65em}
  {}

% Paper-like compact spacing -------------------------------------------------
\titlespacing*{\section}
  {0pt}
  {1.25em plus 0.2em minus 0.1em}
  {0.65em plus 0.1em minus 0.1em}

\titlespacing*{\subsection}
  {0pt}
  {1.0em plus 0.2em minus 0.1em}
  {0.45em plus 0.1em minus 0.1em}

\titlespacing*{\subsubsection}
  {0pt}
  {0.85em plus 0.15em minus 0.1em}
  {0.35em plus 0.1em minus 0.1em}


% NV8: macro bìa phải định nghĩa TRƯỚC \title bên dưới (Overleaf báo Undefined control sequence nếu đặt sau)
\newcommand{\TODO}[1]{\textcolor{red}{\textbf{[TODO: #1]}}}
% Thông tin bìa (NHIEM_VU_8 mục 0.2), người dùng cung cấp; không có người proofreading
\newcommand{\ReportTitle}{DLinear\\[4pt]{\large Linear, DLinear, NLinear tại nghiệm tối ưu: kiến trúc hay cách huấn luyện?}}
\newcommand{\ReportAuthor}{Hoàng Đinh Tuấn Anh}
\newcommand{\ReportDate}{14/10/2026}

% Logo + Title ---------------------------------------------------------------

\setlength{\droptitle}{-7em}

% Title horizontal rule
\newcommand{\TitleRule}{%
  \rule{0.90\textwidth}{0.6pt}
}

% Chiều cao vùng title giữa 2 rule
\newlength{\TitleBlockHeight}
\setlength{\TitleBlockHeight}{2.2cm}

\pretitle{%
  

  % Logo row
  \begin{tabular}{@{}c@{\hspace{1.0em}}c@{\hspace{1.0em}}c@{}}
    \raisebox{-0.35\height}{\includegraphics[height=2.0cm]{Figures/aio2026_logo.pdf}}    
  \end{tabular}

  %\vspace{0.2em}

\begin{center}
  \TitleRule
  \par\nointerlineskip
}

\title{%
  \parbox[c][\TitleBlockHeight][c]{0.90\textwidth}{%
    \centering
    {\fontsize{20pt}{24pt}\selectfont\bfseries \ReportTitle}  % NV8
  }%
}

\posttitle{%
  \par\nointerlineskip
  \TitleRule
  \end{center}
}

\preauthor{%
  \begin{center}
  \fontsize{11pt}{13pt}\selectfont
}

\author{%
  \textbf{\ReportAuthor}\\  % NV8
  \vspace{1em}
  \ReportDate  % NV8 (bỏ dòng Proofreading)
}

\postauthor{%
  \par
  \end{center}
}

\date{}
% Header --------------------------------------------------------------------
\pagestyle{fancy}
\fancyhf{}
\lhead{\href{https://www.facebook.com/aivietnam.edu.vn}{\textcolor{aivnGreen}{\bfseries AI VIET NAM (AIO2026)}}}
\rhead{\href{https://aivietnam.edu.vn/}{\textcolor{aivnGreen}{\bfseries aivietnam.edu.vn}}}
\fancyfoot[r]{\textcolor{aivnGreen}{\bfseries Trang \thepage}}
\fancyfoot[l]{%
  \href{https://zalo.me/84911118758}{%
    \textcolor{aivnGreen}{\bfseries SĐT/Zalo: 0911118758 (Dr. Hồng Phúc)}%
  }%
}
\renewcommand{\headrulewidth}{1.0pt}
\renewcommand{\headrule}{{\color{aivnGreen}\hrule width\headwidth height\headrulewidth \vskip-\headrulewidth}}

\renewcommand{\footrulewidth}{1.0pt}
\renewcommand{\footrule}{{\color{aivnGreen}\hrule width\headwidth height\headrulewidth \vskip-\headrulewidth}}


\fancypagestyle{firstpage}{
    \fancyhf{}
    \fancyfoot[r]{\textcolor{aivnGreen}{\bfseries Trang \thepage}}
    \fancyfoot[l]{%
      \href{https://zalo.me/84911118758}{%
        \textcolor{aivnGreen}{\bfseries SĐT/Zalo: 0911118758 (Dr. Hồng Phúc)}%
      }%
    }
    \renewcommand{\headrulewidth}{0pt}
    \renewcommand{\footrulewidth}{1.0pt}
    \renewcommand{\footrule}{%
        {\color{aivnGreen}\hrule width\headwidth height\footrulewidth}
    }
}



% ---------------------------------------------------------------------------
% NV8: phần thêm cho báo cáo (không đổi kiểu của template)
% ---------------------------------------------------------------------------
\graphicspath{{template/}{figures/}}
% Bảng: chú thích kiểu captionbelowboxaivn của template, thêm bộ đếm riêng để \label, \ref dùng được
\newcounter{bang}
\newcommand{\captab}[1]{\refstepcounter{bang}Bảng~\thebang:~#1}
\newcommand{\RQ}[1]{\textbf{RQ#1}}
\newcommand{\figref}[1]{Hình~\ref{#1}}
\newcommand{\tabref}[1]{Bảng~\ref{#1}}
\newcommand{\secref}[1]{mục~\ref{#1}}
\newcommand{\Weff}{W_{\mathrm{eff}}}
\newcommand{\dmad}{\delta_{\mathrm{MAD}}}
% Tự động sinh bởi scripts/report/make_numbers.py từ results/. KHÔNG SỬA TAY.
%
% Quy ước tên macro: \nm<Nhóm><Tên>, chỉ gồm chữ cái (LaTeX không cho chữ số trong tên lệnh):
%   - mỗi chữ số viết bằng chữ tiếng Anh viết hoa đầu: 1 → One, 96 → NineSix, 720 → SevenTwoZero;
%   - dataset: ETThOne, ETThTwo, ETTmOne; horizon: H + số (HNineSix); mô hình: Linear, DLinear, NLinear;
%   - mục tiêu: MSE, MAE, Huber; metric báo cáo: Mse (MSE test ở λ*_MSE), Mae (MAE test ở λ*_MAE);
%   - nhóm: BTwo (RQ1, bảng B2), BThree (B3, nghiệm tối ưu), Out (outlier), Mag (độ lớn outlier), Delta (δ của Huber),
%     BOne, BFour, CThree (SGD với lr gốc), Tuned (SGD đã tinh chỉnh lr, lưới đã mở rộng), EThTwo (đoạn cảm biến bị kẹt),
%     Drop (code gốc với drop_last), Ridge (họ ridge trên ETTh2), Stuck (vị trí đoạn hằng), DeltaH (δ theo bước h);
%   - hậu tố Min/Max/Range/Count/Pct là khoảng hoặc số đếm; số đếm không có đơn vị.
% Số thập phân dùng dấu phẩy, bọc trong \ensuremath nên dùng được cả trong văn bản lẫn công thức.

\newcommand{\nmReproGridHours}{\ensuremath{14{,}6}}
\newcommand{\nmReproGridHoursETThOne}{\ensuremath{3{,}5}}
\newcommand{\nmReproGridHoursETThTwo}{\ensuremath{3{,}9}}
\newcommand{\nmReproGridHoursETTmOne}{\ensuremath{7{,}3}}
\newcommand{\nmReproCheckCellMinutes}{14}
\newcommand{\nmReproHuberMadHours}{\ensuremath{0{,}6}}
\newcommand{\nmReproSgdCpuHours}{\ensuremath{5{,}5}}
\newcommand{\nmReproSgdLrWallHours}{\ensuremath{2{,}3}}
\newcommand{\nmReproHead}{\texttt{d3c31e2}}
\newcommand{\nmBTwoRankZ}{336}
\newcommand{\nmBTwoNcolZ}{672}
\newcommand{\nmBTwoDLinZeroRelDWMax}{\ensuremath{5\cdot 10^{-13}}}
\newcommand{\nmBTwoDLinZeroDyhatMax}{\ensuremath{2\cdot 10^{-13}}}
\newcommand{\nmBTwoDLinZeroGrad}{\ensuremath{6\cdot 10^{-15}}}
\newcommand{\nmBTwoDLinRelDWLamFourKFive}{\ensuremath{0{,}12}}
\newcommand{\nmBTwoDLinRmsLamFourKFive}{\ensuremath{0{,}017}}
\newcommand{\nmBTwoDLinRmsBestLamFourKFive}{\ensuremath{0{,}015}}
\newcommand{\nmBTwoDLinRelDWLamFourKOneFive}{\ensuremath{0{,}19}}
\newcommand{\nmBTwoDLinRmsLamFourKOneFive}{\ensuremath{0{,}018}}
\newcommand{\nmBTwoDLinRmsBestLamFourKOneFive}{\ensuremath{0{,}018}}
\newcommand{\nmBTwoDLinRelDWLamFourKTwoFive}{\ensuremath{0{,}28}}
\newcommand{\nmBTwoDLinRmsLamFourKTwoFive}{\ensuremath{0{,}022}}
\newcommand{\nmBTwoDLinRmsBestLamFourKTwoFive}{\ensuremath{0{,}021}}
\newcommand{\nmBTwoDLinRelDWLamFourKFourNine}{\ensuremath{0{,}37}}
\newcommand{\nmBTwoDLinRmsLamFourKFourNine}{\ensuremath{0{,}029}}
\newcommand{\nmBTwoDLinRmsBestLamFourKFourNine}{\ensuremath{0{,}027}}
\newcommand{\nmBTwoDLinBestOverSameMedianPct}{91}
\newcommand{\nmBTwoNLinDWF}{\ensuremath{0{,}0694}}
\newcommand{\nmBTwoNLinRmsZero}{\ensuremath{0{,}0514}}
\newcommand{\nmBTwoNLinDSSE}{\ensuremath{17002{,}8}}
\newcommand{\nmBTwoNLinRmsLamFour}{\ensuremath{0{,}069}}
\newcommand{\nmBTwoNLinRmsLamEight}{\ensuremath{0{,}973}}
\newcommand{\nmBTwoLinRmseTest}{\ensuremath{0{,}608}}
\newcommand{\nmBThreeDLinMinusLinAbsMax}{\ensuremath{0{,}0006}}
\newcommand{\nmBThreeDLinMinusLinAllAbsMax}{\ensuremath{0{,}0016}}
\newcommand{\nmBThreeDLinMinusLinAllCount}{72}
\newcommand{\nmBThreeNLinMinusLinETThOneRange}{\ensuremath{-0{,}0007}--\ensuremath{-0{,}0371}}
\newcommand{\nmBThreeNLinMinusLinETThTwoRange}{\ensuremath{-0{,}0292}--\ensuremath{-0{,}3480}}
\newcommand{\nmBThreeNLinMinusLinETTmOneRange}{\ensuremath{0{,}0012}--\ensuremath{0{,}0030}}
\newcommand{\nmCThreeCfAbsMax}{\ensuremath{0{,}0006}}
\newcommand{\nmBThreeETThOneHNineSixLinearMSEMse}{\ensuremath{0{,}3702}}
\newcommand{\nmBThreeETThOneHNineSixLinearMSEMae}{\ensuremath{0{,}3915}}
\newcommand{\nmBThreeETThOneHNineSixLinearMAEMse}{\ensuremath{0{,}3653}}
\newcommand{\nmBThreeETThOneHNineSixLinearMAEMae}{\ensuremath{0{,}3833}}
\newcommand{\nmBThreeETThOneHNineSixLinearHuberMse}{\ensuremath{0{,}3658}}
\newcommand{\nmBThreeETThOneHNineSixLinearHuberMae}{\ensuremath{0{,}3852}}
\newcommand{\nmBThreeETThOneHNineSixLinearMSELam}{\ensuremath{0}}
\newcommand{\nmDeltaETThOneHNineSixLinear}{\ensuremath{0{,}538}}
\newcommand{\nmBThreeETThOneHNineSixDLinearMSEMse}{\ensuremath{0{,}3698}}
\newcommand{\nmBThreeETThOneHNineSixDLinearMSEMae}{\ensuremath{0{,}3914}}
\newcommand{\nmBThreeETThOneHNineSixDLinearMAEMse}{\ensuremath{0{,}3647}}
\newcommand{\nmBThreeETThOneHNineSixDLinearMAEMae}{\ensuremath{0{,}3832}}
\newcommand{\nmBThreeETThOneHNineSixDLinearHuberMse}{\ensuremath{0{,}3652}}
\newcommand{\nmBThreeETThOneHNineSixDLinearHuberMae}{\ensuremath{0{,}3851}}
\newcommand{\nmBThreeETThOneHNineSixDLinearMSELam}{\ensuremath{446{,}7}}
\newcommand{\nmDeltaETThOneHNineSixDLinear}{\ensuremath{0{,}538}}
\newcommand{\nmBThreeETThOneHNineSixNLinearMSEMse}{\ensuremath{0{,}3695}}
\newcommand{\nmBThreeETThOneHNineSixNLinearMSEMae}{\ensuremath{0{,}3913}}
\newcommand{\nmBThreeETThOneHNineSixNLinearMAEMse}{\ensuremath{0{,}3638}}
\newcommand{\nmBThreeETThOneHNineSixNLinearMAEMae}{\ensuremath{0{,}3839}}
\newcommand{\nmBThreeETThOneHNineSixNLinearHuberMse}{\ensuremath{0{,}3656}}
\newcommand{\nmBThreeETThOneHNineSixNLinearHuberMae}{\ensuremath{0{,}3860}}
\newcommand{\nmBThreeETThOneHNineSixNLinearMSELam}{\ensuremath{1413}}
\newcommand{\nmDeltaETThOneHNineSixNLinear}{\ensuremath{0{,}538}}
\newcommand{\nmBThreeETThOneHOneNineTwoLinearMSEMse}{\ensuremath{0{,}4040}}
\newcommand{\nmBThreeETThOneHOneNineTwoLinearMSEMae}{\ensuremath{0{,}4126}}
\newcommand{\nmBThreeETThOneHOneNineTwoLinearMAEMse}{\ensuremath{0{,}4031}}
\newcommand{\nmBThreeETThOneHOneNineTwoLinearMAEMae}{\ensuremath{0{,}4074}}
\newcommand{\nmBThreeETThOneHOneNineTwoLinearHuberMse}{\ensuremath{0{,}4025}}
\newcommand{\nmBThreeETThOneHOneNineTwoLinearHuberMae}{\ensuremath{0{,}4082}}
\newcommand{\nmBThreeETThOneHOneNineTwoLinearMSELam}{\ensuremath{223{,}9}}
\newcommand{\nmDeltaETThOneHOneNineTwoLinear}{\ensuremath{0{,}582}}
\newcommand{\nmBThreeETThOneHOneNineTwoDLinearMSEMse}{\ensuremath{0{,}4036}}
\newcommand{\nmBThreeETThOneHOneNineTwoDLinearMSEMae}{\ensuremath{0{,}4125}}
\newcommand{\nmBThreeETThOneHOneNineTwoDLinearMAEMse}{\ensuremath{0{,}4014}}
\newcommand{\nmBThreeETThOneHOneNineTwoDLinearMAEMae}{\ensuremath{0{,}4074}}
\newcommand{\nmBThreeETThOneHOneNineTwoDLinearHuberMse}{\ensuremath{0{,}4017}}
\newcommand{\nmBThreeETThOneHOneNineTwoDLinearHuberMae}{\ensuremath{0{,}4080}}
\newcommand{\nmBThreeETThOneHOneNineTwoDLinearMSELam}{\ensuremath{1000}}
\newcommand{\nmDeltaETThOneHOneNineTwoDLinear}{\ensuremath{0{,}582}}
\newcommand{\nmBThreeETThOneHOneNineTwoNLinearMSEMse}{\ensuremath{0{,}4027}}
\newcommand{\nmBThreeETThOneHOneNineTwoNLinearMSEMae}{\ensuremath{0{,}4110}}
\newcommand{\nmBThreeETThOneHOneNineTwoNLinearMAEMse}{\ensuremath{0{,}4026}}
\newcommand{\nmBThreeETThOneHOneNineTwoNLinearMAEMae}{\ensuremath{0{,}4075}}
\newcommand{\nmBThreeETThOneHOneNineTwoNLinearHuberMse}{\ensuremath{0{,}4024}}
\newcommand{\nmBThreeETThOneHOneNineTwoNLinearHuberMae}{\ensuremath{0{,}4084}}
\newcommand{\nmBThreeETThOneHOneNineTwoNLinearMSELam}{\ensuremath{2512}}
\newcommand{\nmDeltaETThOneHOneNineTwoNLinear}{\ensuremath{0{,}587}}
\newcommand{\nmBThreeETThOneHThreeThreeSixLinearMSEMse}{\ensuremath{0{,}4324}}
\newcommand{\nmBThreeETThOneHThreeThreeSixLinearMSEMae}{\ensuremath{0{,}4340}}
\newcommand{\nmBThreeETThOneHThreeThreeSixLinearMAEMse}{\ensuremath{0{,}4302}}
\newcommand{\nmBThreeETThOneHThreeThreeSixLinearMAEMae}{\ensuremath{0{,}4275}}
\newcommand{\nmBThreeETThOneHThreeThreeSixLinearHuberMse}{\ensuremath{0{,}4296}}
\newcommand{\nmBThreeETThOneHThreeThreeSixLinearHuberMae}{\ensuremath{0{,}4285}}
\newcommand{\nmBThreeETThOneHThreeThreeSixLinearMSELam}{\ensuremath{1995}}
\newcommand{\nmDeltaETThOneHThreeThreeSixLinear}{\ensuremath{0{,}621}}
\newcommand{\nmBThreeETThOneHThreeThreeSixDLinearMSEMse}{\ensuremath{0{,}4317}}
\newcommand{\nmBThreeETThOneHThreeThreeSixDLinearMSEMae}{\ensuremath{0{,}4340}}
\newcommand{\nmBThreeETThOneHThreeThreeSixDLinearMAEMse}{\ensuremath{0{,}4296}}
\newcommand{\nmBThreeETThOneHThreeThreeSixDLinearMAEMae}{\ensuremath{0{,}4272}}
\newcommand{\nmBThreeETThOneHThreeThreeSixDLinearHuberMse}{\ensuremath{0{,}4290}}
\newcommand{\nmBThreeETThOneHThreeThreeSixDLinearHuberMae}{\ensuremath{0{,}4285}}
\newcommand{\nmBThreeETThOneHThreeThreeSixDLinearMSELam}{\ensuremath{4467}}
\newcommand{\nmDeltaETThOneHThreeThreeSixDLinear}{\ensuremath{0{,}621}}
\newcommand{\nmBThreeETThOneHThreeThreeSixNLinearMSEMse}{\ensuremath{0{,}4261}}
\newcommand{\nmBThreeETThOneHThreeThreeSixNLinearMSEMae}{\ensuremath{0{,}4260}}
\newcommand{\nmBThreeETThOneHThreeThreeSixNLinearMAEMse}{\ensuremath{0{,}4290}}
\newcommand{\nmBThreeETThOneHThreeThreeSixNLinearMAEMae}{\ensuremath{0{,}4237}}
\newcommand{\nmBThreeETThOneHThreeThreeSixNLinearHuberMse}{\ensuremath{0{,}4307}}
\newcommand{\nmBThreeETThOneHThreeThreeSixNLinearHuberMae}{\ensuremath{0{,}4246}}
\newcommand{\nmBThreeETThOneHThreeThreeSixNLinearMSELam}{\ensuremath{8913}}
\newcommand{\nmDeltaETThOneHThreeThreeSixNLinear}{\ensuremath{0{,}629}}
\newcommand{\nmBThreeETThOneHSevenTwoZeroLinearMSEMse}{\ensuremath{0{,}4701}}
\newcommand{\nmBThreeETThOneHSevenTwoZeroLinearMSEMae}{\ensuremath{0{,}4878}}
\newcommand{\nmBThreeETThOneHSevenTwoZeroLinearMAEMse}{\ensuremath{0{,}4553}}
\newcommand{\nmBThreeETThOneHSevenTwoZeroLinearMAEMae}{\ensuremath{0{,}4747}}
\newcommand{\nmBThreeETThOneHSevenTwoZeroLinearHuberMse}{\ensuremath{0{,}4624}}
\newcommand{\nmBThreeETThOneHSevenTwoZeroLinearHuberMae}{\ensuremath{0{,}4767}}
\newcommand{\nmBThreeETThOneHSevenTwoZeroLinearMSELam}{\ensuremath{1{,}778\cdot 10^{4}}}
\newcommand{\nmDeltaETThOneHSevenTwoZeroLinear}{\ensuremath{0{,}709}}
\newcommand{\nmBThreeETThOneHSevenTwoZeroDLinearMSEMse}{\ensuremath{0{,}4697}}
\newcommand{\nmBThreeETThOneHSevenTwoZeroDLinearMSEMae}{\ensuremath{0{,}4879}}
\newcommand{\nmBThreeETThOneHSevenTwoZeroDLinearMAEMse}{\ensuremath{0{,}4549}}
\newcommand{\nmBThreeETThOneHSevenTwoZeroDLinearMAEMae}{\ensuremath{0{,}4746}}
\newcommand{\nmBThreeETThOneHSevenTwoZeroDLinearHuberMse}{\ensuremath{0{,}4620}}
\newcommand{\nmBThreeETThOneHSevenTwoZeroDLinearHuberMae}{\ensuremath{0{,}4766}}
\newcommand{\nmBThreeETThOneHSevenTwoZeroDLinearMSELam}{\ensuremath{1{,}778\cdot 10^{4}}}
\newcommand{\nmDeltaETThOneHSevenTwoZeroDLinear}{\ensuremath{0{,}709}}
\newcommand{\nmBThreeETThOneHSevenTwoZeroNLinearMSEMse}{\ensuremath{0{,}4329}}
\newcommand{\nmBThreeETThOneHSevenTwoZeroNLinearMSEMae}{\ensuremath{0{,}4505}}
\newcommand{\nmBThreeETThOneHSevenTwoZeroNLinearMAEMse}{\ensuremath{0{,}4282}}
\newcommand{\nmBThreeETThOneHSevenTwoZeroNLinearMAEMae}{\ensuremath{0{,}4454}}
\newcommand{\nmBThreeETThOneHSevenTwoZeroNLinearHuberMse}{\ensuremath{0{,}4329}}
\newcommand{\nmBThreeETThOneHSevenTwoZeroNLinearHuberMae}{\ensuremath{0{,}4468}}
\newcommand{\nmBThreeETThOneHSevenTwoZeroNLinearMSELam}{\ensuremath{1{,}995\cdot 10^{4}}}
\newcommand{\nmDeltaETThOneHSevenTwoZeroNLinear}{\ensuremath{0{,}731}}
\newcommand{\nmBThreeETThTwoHNineSixLinearMSEMse}{\ensuremath{0{,}3009}}
\newcommand{\nmBThreeETThTwoHNineSixLinearMSEMae}{\ensuremath{0{,}3649}}
\newcommand{\nmBThreeETThTwoHNineSixLinearMAEMse}{\ensuremath{0{,}2779}}
\newcommand{\nmBThreeETThTwoHNineSixLinearMAEMae}{\ensuremath{0{,}3334}}
\newcommand{\nmBThreeETThTwoHNineSixLinearHuberMse}{\ensuremath{0{,}2818}}
\newcommand{\nmBThreeETThTwoHNineSixLinearHuberMae}{\ensuremath{0{,}3410}}
\newcommand{\nmBThreeETThTwoHNineSixLinearMSELam}{\ensuremath{1585}}
\newcommand{\nmDeltaETThTwoHNineSixLinear}{\ensuremath{0{,}398}}
\newcommand{\nmBThreeETThTwoHNineSixDLinearMSEMse}{\ensuremath{0{,}3008}}
\newcommand{\nmBThreeETThTwoHNineSixDLinearMSEMae}{\ensuremath{0{,}3649}}
\newcommand{\nmBThreeETThTwoHNineSixDLinearMAEMse}{\ensuremath{0{,}2778}}
\newcommand{\nmBThreeETThTwoHNineSixDLinearMAEMae}{\ensuremath{0{,}3333}}
\newcommand{\nmBThreeETThTwoHNineSixDLinearHuberMse}{\ensuremath{0{,}2817}}
\newcommand{\nmBThreeETThTwoHNineSixDLinearHuberMae}{\ensuremath{0{,}3410}}
\newcommand{\nmBThreeETThTwoHNineSixDLinearMSELam}{\ensuremath{1995}}
\newcommand{\nmDeltaETThTwoHNineSixDLinear}{\ensuremath{0{,}398}}
\newcommand{\nmBThreeETThTwoHNineSixNLinearMSEMse}{\ensuremath{0{,}2717}}
\newcommand{\nmBThreeETThTwoHNineSixNLinearMSEMae}{\ensuremath{0{,}3342}}
\newcommand{\nmBThreeETThTwoHNineSixNLinearMAEMse}{\ensuremath{0{,}2764}}
\newcommand{\nmBThreeETThTwoHNineSixNLinearMAEMae}{\ensuremath{0{,}3307}}
\newcommand{\nmBThreeETThTwoHNineSixNLinearHuberMse}{\ensuremath{0{,}2756}}
\newcommand{\nmBThreeETThTwoHNineSixNLinearHuberMae}{\ensuremath{0{,}3306}}
\newcommand{\nmBThreeETThTwoHNineSixNLinearMSELam}{\ensuremath{1778}}
\newcommand{\nmDeltaETThTwoHNineSixNLinear}{\ensuremath{0{,}372}}
\newcommand{\nmBThreeETThTwoHOneNineTwoLinearMSEMse}{\ensuremath{0{,}3966}}
\newcommand{\nmBThreeETThTwoHOneNineTwoLinearMSEMae}{\ensuremath{0{,}4273}}
\newcommand{\nmBThreeETThTwoHOneNineTwoLinearMAEMse}{\ensuremath{0{,}3531}}
\newcommand{\nmBThreeETThTwoHOneNineTwoLinearMAEMae}{\ensuremath{0{,}3855}}
\newcommand{\nmBThreeETThTwoHOneNineTwoLinearHuberMse}{\ensuremath{0{,}3600}}
\newcommand{\nmBThreeETThTwoHOneNineTwoLinearHuberMae}{\ensuremath{0{,}3945}}
\newcommand{\nmBThreeETThTwoHOneNineTwoLinearMSELam}{\ensuremath{1122}}
\newcommand{\nmDeltaETThTwoHOneNineTwoLinear}{\ensuremath{0{,}453}}
\newcommand{\nmBThreeETThTwoHOneNineTwoDLinearMSEMse}{\ensuremath{0{,}3966}}
\newcommand{\nmBThreeETThTwoHOneNineTwoDLinearMSEMae}{\ensuremath{0{,}4273}}
\newcommand{\nmBThreeETThTwoHOneNineTwoDLinearMAEMse}{\ensuremath{0{,}3531}}
\newcommand{\nmBThreeETThTwoHOneNineTwoDLinearMAEMae}{\ensuremath{0{,}3855}}
\newcommand{\nmBThreeETThTwoHOneNineTwoDLinearHuberMse}{\ensuremath{0{,}3600}}
\newcommand{\nmBThreeETThTwoHOneNineTwoDLinearHuberMae}{\ensuremath{0{,}3945}}
\newcommand{\nmBThreeETThTwoHOneNineTwoDLinearMSELam}{\ensuremath{1259}}
\newcommand{\nmDeltaETThTwoHOneNineTwoDLinear}{\ensuremath{0{,}453}}
\newcommand{\nmBThreeETThTwoHOneNineTwoNLinearMSEMse}{\ensuremath{0{,}3336}}
\newcommand{\nmBThreeETThTwoHOneNineTwoNLinearMSEMae}{\ensuremath{0{,}3754}}
\newcommand{\nmBThreeETThTwoHOneNineTwoNLinearMAEMse}{\ensuremath{0{,}3447}}
\newcommand{\nmBThreeETThTwoHOneNineTwoNLinearMAEMae}{\ensuremath{0{,}3756}}
\newcommand{\nmBThreeETThTwoHOneNineTwoNLinearHuberMse}{\ensuremath{0{,}3427}}
\newcommand{\nmBThreeETThTwoHOneNineTwoNLinearHuberMae}{\ensuremath{0{,}3744}}
\newcommand{\nmBThreeETThTwoHOneNineTwoNLinearMSELam}{\ensuremath{1585}}
\newcommand{\nmDeltaETThTwoHOneNineTwoNLinear}{\ensuremath{0{,}421}}
\newcommand{\nmBThreeETThTwoHThreeThreeSixLinearMSEMse}{\ensuremath{0{,}4853}}
\newcommand{\nmBThreeETThTwoHThreeThreeSixLinearMSEMae}{\ensuremath{0{,}4835}}
\newcommand{\nmBThreeETThTwoHThreeThreeSixLinearMAEMse}{\ensuremath{0{,}4106}}
\newcommand{\nmBThreeETThTwoHThreeThreeSixLinearMAEMae}{\ensuremath{0{,}4296}}
\newcommand{\nmBThreeETThTwoHThreeThreeSixLinearHuberMse}{\ensuremath{0{,}4220}}
\newcommand{\nmBThreeETThTwoHThreeThreeSixLinearHuberMae}{\ensuremath{0{,}4400}}
\newcommand{\nmBThreeETThTwoHThreeThreeSixLinearMSELam}{\ensuremath{1000}}
\newcommand{\nmDeltaETThTwoHThreeThreeSixLinear}{\ensuremath{0{,}505}}
\newcommand{\nmBThreeETThTwoHThreeThreeSixDLinearMSEMse}{\ensuremath{0{,}4853}}
\newcommand{\nmBThreeETThTwoHThreeThreeSixDLinearMSEMae}{\ensuremath{0{,}4835}}
\newcommand{\nmBThreeETThTwoHThreeThreeSixDLinearMAEMse}{\ensuremath{0{,}4106}}
\newcommand{\nmBThreeETThTwoHThreeThreeSixDLinearMAEMae}{\ensuremath{0{,}4296}}
\newcommand{\nmBThreeETThTwoHThreeThreeSixDLinearHuberMse}{\ensuremath{0{,}4220}}
\newcommand{\nmBThreeETThTwoHThreeThreeSixDLinearHuberMae}{\ensuremath{0{,}4400}}
\newcommand{\nmBThreeETThTwoHThreeThreeSixDLinearMSELam}{\ensuremath{1122}}
\newcommand{\nmDeltaETThTwoHThreeThreeSixDLinear}{\ensuremath{0{,}505}}
\newcommand{\nmBThreeETThTwoHThreeThreeSixNLinearMSEMse}{\ensuremath{0{,}3581}}
\newcommand{\nmBThreeETThTwoHThreeThreeSixNLinearMSEMae}{\ensuremath{0{,}3997}}
\newcommand{\nmBThreeETThTwoHThreeThreeSixNLinearMAEMse}{\ensuremath{0{,}3760}}
\newcommand{\nmBThreeETThTwoHThreeThreeSixNLinearMAEMae}{\ensuremath{0{,}4028}}
\newcommand{\nmBThreeETThTwoHThreeThreeSixNLinearHuberMse}{\ensuremath{0{,}3727}}
\newcommand{\nmBThreeETThTwoHThreeThreeSixNLinearHuberMae}{\ensuremath{0{,}4008}}
\newcommand{\nmBThreeETThTwoHThreeThreeSixNLinearMSELam}{\ensuremath{1585}}
\newcommand{\nmDeltaETThTwoHThreeThreeSixNLinear}{\ensuremath{0{,}470}}
\newcommand{\nmBThreeETThTwoHSevenTwoZeroLinearMSEMse}{\ensuremath{0{,}7404}}
\newcommand{\nmBThreeETThTwoHSevenTwoZeroLinearMSEMae}{\ensuremath{0{,}6115}}
\newcommand{\nmBThreeETThTwoHSevenTwoZeroLinearMAEMse}{\ensuremath{0{,}5893}}
\newcommand{\nmBThreeETThTwoHSevenTwoZeroLinearMAEMae}{\ensuremath{0{,}5335}}
\newcommand{\nmBThreeETThTwoHSevenTwoZeroLinearHuberMse}{\ensuremath{0{,}6277}}
\newcommand{\nmBThreeETThTwoHSevenTwoZeroLinearHuberMae}{\ensuremath{0{,}5541}}
\newcommand{\nmBThreeETThTwoHSevenTwoZeroLinearMSELam}{\ensuremath{1000}}
\newcommand{\nmDeltaETThTwoHSevenTwoZeroLinear}{\ensuremath{0{,}610}}
\newcommand{\nmBThreeETThTwoHSevenTwoZeroDLinearMSEMse}{\ensuremath{0{,}7403}}
\newcommand{\nmBThreeETThTwoHSevenTwoZeroDLinearMSEMae}{\ensuremath{0{,}6115}}
\newcommand{\nmBThreeETThTwoHSevenTwoZeroDLinearMAEMse}{\ensuremath{0{,}5892}}
\newcommand{\nmBThreeETThTwoHSevenTwoZeroDLinearMAEMae}{\ensuremath{0{,}5335}}
\newcommand{\nmBThreeETThTwoHSevenTwoZeroDLinearHuberMse}{\ensuremath{0{,}6276}}
\newcommand{\nmBThreeETThTwoHSevenTwoZeroDLinearHuberMae}{\ensuremath{0{,}5541}}
\newcommand{\nmBThreeETThTwoHSevenTwoZeroDLinearMSELam}{\ensuremath{1122}}
\newcommand{\nmDeltaETThTwoHSevenTwoZeroDLinear}{\ensuremath{0{,}610}}
\newcommand{\nmBThreeETThTwoHSevenTwoZeroNLinearMSEMse}{\ensuremath{0{,}3924}}
\newcommand{\nmBThreeETThTwoHSevenTwoZeroNLinearMSEMae}{\ensuremath{0{,}4341}}
\newcommand{\nmBThreeETThTwoHSevenTwoZeroNLinearMAEMse}{\ensuremath{0{,}3926}}
\newcommand{\nmBThreeETThTwoHSevenTwoZeroNLinearMAEMae}{\ensuremath{0{,}4242}}
\newcommand{\nmBThreeETThTwoHSevenTwoZeroNLinearHuberMse}{\ensuremath{0{,}3912}}
\newcommand{\nmBThreeETThTwoHSevenTwoZeroNLinearHuberMae}{\ensuremath{0{,}4251}}
\newcommand{\nmBThreeETThTwoHSevenTwoZeroNLinearMSELam}{\ensuremath{1995}}
\newcommand{\nmDeltaETThTwoHSevenTwoZeroNLinear}{\ensuremath{0{,}587}}
\newcommand{\nmBThreeETTmOneHNineSixLinearMSEMse}{\ensuremath{0{,}2992}}
\newcommand{\nmBThreeETTmOneHNineSixLinearMSEMae}{\ensuremath{0{,}3427}}
\newcommand{\nmBThreeETTmOneHNineSixLinearMAEMse}{\ensuremath{0{,}2929}}
\newcommand{\nmBThreeETTmOneHNineSixLinearMAEMae}{\ensuremath{0{,}3331}}
\newcommand{\nmBThreeETTmOneHNineSixLinearHuberMse}{\ensuremath{0{,}2953}}
\newcommand{\nmBThreeETTmOneHNineSixLinearHuberMae}{\ensuremath{0{,}3362}}
\newcommand{\nmBThreeETTmOneHNineSixLinearMSELam}{\ensuremath{631}}
\newcommand{\nmDeltaETTmOneHNineSixLinear}{\ensuremath{0{,}455}}
\newcommand{\nmBThreeETTmOneHNineSixDLinearMSEMse}{\ensuremath{0{,}2992}}
\newcommand{\nmBThreeETTmOneHNineSixDLinearMSEMae}{\ensuremath{0{,}3427}}
\newcommand{\nmBThreeETTmOneHNineSixDLinearMAEMse}{\ensuremath{0{,}2929}}
\newcommand{\nmBThreeETTmOneHNineSixDLinearMAEMae}{\ensuremath{0{,}3331}}
\newcommand{\nmBThreeETTmOneHNineSixDLinearHuberMse}{\ensuremath{0{,}2952}}
\newcommand{\nmBThreeETTmOneHNineSixDLinearHuberMae}{\ensuremath{0{,}3362}}
\newcommand{\nmBThreeETTmOneHNineSixDLinearMSELam}{\ensuremath{1413}}
\newcommand{\nmDeltaETTmOneHNineSixDLinear}{\ensuremath{0{,}455}}
\newcommand{\nmBThreeETTmOneHNineSixNLinearMSEMse}{\ensuremath{0{,}3004}}
\newcommand{\nmBThreeETTmOneHNineSixNLinearMSEMae}{\ensuremath{0{,}3426}}
\newcommand{\nmBThreeETTmOneHNineSixNLinearMAEMse}{\ensuremath{0{,}2937}}
\newcommand{\nmBThreeETTmOneHNineSixNLinearMAEMae}{\ensuremath{0{,}3337}}
\newcommand{\nmBThreeETTmOneHNineSixNLinearHuberMse}{\ensuremath{0{,}2964}}
\newcommand{\nmBThreeETTmOneHNineSixNLinearHuberMae}{\ensuremath{0{,}3367}}
\newcommand{\nmBThreeETTmOneHNineSixNLinearMSELam}{\ensuremath{5012}}
\newcommand{\nmDeltaETTmOneHNineSixNLinear}{\ensuremath{0{,}453}}
\newcommand{\nmBThreeETTmOneHOneNineTwoLinearMSEMse}{\ensuremath{0{,}3340}}
\newcommand{\nmBThreeETTmOneHOneNineTwoLinearMSEMae}{\ensuremath{0{,}3637}}
\newcommand{\nmBThreeETTmOneHOneNineTwoLinearMAEMse}{\ensuremath{0{,}3322}}
\newcommand{\nmBThreeETTmOneHOneNineTwoLinearMAEMae}{\ensuremath{0{,}3553}}
\newcommand{\nmBThreeETTmOneHOneNineTwoLinearHuberMse}{\ensuremath{0{,}3327}}
\newcommand{\nmBThreeETTmOneHOneNineTwoLinearHuberMae}{\ensuremath{0{,}3578}}
\newcommand{\nmBThreeETTmOneHOneNineTwoLinearMSELam}{\ensuremath{354{,}8}}
\newcommand{\nmDeltaETTmOneHOneNineTwoLinear}{\ensuremath{0{,}497}}
\newcommand{\nmBThreeETTmOneHOneNineTwoDLinearMSEMse}{\ensuremath{0{,}3339}}
\newcommand{\nmBThreeETTmOneHOneNineTwoDLinearMSEMae}{\ensuremath{0{,}3637}}
\newcommand{\nmBThreeETTmOneHOneNineTwoDLinearMAEMse}{\ensuremath{0{,}3319}}
\newcommand{\nmBThreeETTmOneHOneNineTwoDLinearMAEMae}{\ensuremath{0{,}3553}}
\newcommand{\nmBThreeETTmOneHOneNineTwoDLinearHuberMse}{\ensuremath{0{,}3326}}
\newcommand{\nmBThreeETTmOneHOneNineTwoDLinearHuberMae}{\ensuremath{0{,}3578}}
\newcommand{\nmBThreeETTmOneHOneNineTwoDLinearMSELam}{\ensuremath{1122}}
\newcommand{\nmDeltaETTmOneHOneNineTwoDLinear}{\ensuremath{0{,}497}}
\newcommand{\nmBThreeETTmOneHOneNineTwoNLinearMSEMse}{\ensuremath{0{,}3356}}
\newcommand{\nmBThreeETTmOneHOneNineTwoNLinearMSEMae}{\ensuremath{0{,}3634}}
\newcommand{\nmBThreeETTmOneHOneNineTwoNLinearMAEMse}{\ensuremath{0{,}3335}}
\newcommand{\nmBThreeETTmOneHOneNineTwoNLinearMAEMae}{\ensuremath{0{,}3563}}
\newcommand{\nmBThreeETTmOneHOneNineTwoNLinearHuberMse}{\ensuremath{0{,}3345}}
\newcommand{\nmBThreeETTmOneHOneNineTwoNLinearHuberMae}{\ensuremath{0{,}3585}}
\newcommand{\nmBThreeETTmOneHOneNineTwoNLinearMSELam}{\ensuremath{5623}}
\newcommand{\nmDeltaETTmOneHOneNineTwoNLinear}{\ensuremath{0{,}493}}
\newcommand{\nmBThreeETTmOneHThreeThreeSixLinearMSEMse}{\ensuremath{0{,}3683}}
\newcommand{\nmBThreeETTmOneHThreeThreeSixLinearMSEMae}{\ensuremath{0{,}3843}}
\newcommand{\nmBThreeETTmOneHThreeThreeSixLinearMAEMse}{\ensuremath{0{,}3658}}
\newcommand{\nmBThreeETTmOneHThreeThreeSixLinearMAEMae}{\ensuremath{0{,}3759}}
\newcommand{\nmBThreeETTmOneHThreeThreeSixLinearHuberMse}{\ensuremath{0{,}3661}}
\newcommand{\nmBThreeETTmOneHThreeThreeSixLinearHuberMae}{\ensuremath{0{,}3782}}
\newcommand{\nmBThreeETTmOneHThreeThreeSixLinearMSELam}{\ensuremath{446{,}7}}
\newcommand{\nmDeltaETTmOneHThreeThreeSixLinear}{\ensuremath{0{,}536}}
\newcommand{\nmBThreeETTmOneHThreeThreeSixDLinearMSEMse}{\ensuremath{0{,}3683}}
\newcommand{\nmBThreeETTmOneHThreeThreeSixDLinearMSEMae}{\ensuremath{0{,}3843}}
\newcommand{\nmBThreeETTmOneHThreeThreeSixDLinearMAEMse}{\ensuremath{0{,}3655}}
\newcommand{\nmBThreeETTmOneHThreeThreeSixDLinearMAEMae}{\ensuremath{0{,}3759}}
\newcommand{\nmBThreeETTmOneHThreeThreeSixDLinearHuberMse}{\ensuremath{0{,}3660}}
\newcommand{\nmBThreeETTmOneHThreeThreeSixDLinearHuberMae}{\ensuremath{0{,}3782}}
\newcommand{\nmBThreeETTmOneHThreeThreeSixDLinearMSELam}{\ensuremath{1259}}
\newcommand{\nmDeltaETTmOneHThreeThreeSixDLinear}{\ensuremath{0{,}536}}
\newcommand{\nmBThreeETTmOneHThreeThreeSixNLinearMSEMse}{\ensuremath{0{,}3703}}
\newcommand{\nmBThreeETTmOneHThreeThreeSixNLinearMSEMae}{\ensuremath{0{,}3834}}
\newcommand{\nmBThreeETTmOneHThreeThreeSixNLinearMAEMse}{\ensuremath{0{,}3685}}
\newcommand{\nmBThreeETTmOneHThreeThreeSixNLinearMAEMae}{\ensuremath{0{,}3770}}
\newcommand{\nmBThreeETTmOneHThreeThreeSixNLinearHuberMse}{\ensuremath{0{,}3691}}
\newcommand{\nmBThreeETTmOneHThreeThreeSixNLinearHuberMae}{\ensuremath{0{,}3790}}
\newcommand{\nmBThreeETTmOneHThreeThreeSixNLinearMSELam}{\ensuremath{6310}}
\newcommand{\nmDeltaETTmOneHThreeThreeSixNLinear}{\ensuremath{0{,}534}}
\newcommand{\nmBThreeETTmOneHSevenTwoZeroLinearMSEMse}{\ensuremath{0{,}4232}}
\newcommand{\nmBThreeETTmOneHSevenTwoZeroLinearMSEMae}{\ensuremath{0{,}4181}}
\newcommand{\nmBThreeETTmOneHSevenTwoZeroLinearMAEMse}{\ensuremath{0{,}4216}}
\newcommand{\nmBThreeETTmOneHSevenTwoZeroLinearMAEMae}{\ensuremath{0{,}4101}}
\newcommand{\nmBThreeETTmOneHSevenTwoZeroLinearHuberMse}{\ensuremath{0{,}4210}}
\newcommand{\nmBThreeETTmOneHSevenTwoZeroLinearHuberMae}{\ensuremath{0{,}4120}}
\newcommand{\nmBThreeETTmOneHSevenTwoZeroLinearMSELam}{\ensuremath{1259}}
\newcommand{\nmDeltaETTmOneHSevenTwoZeroLinear}{\ensuremath{0{,}589}}
\newcommand{\nmBThreeETTmOneHSevenTwoZeroDLinearMSEMse}{\ensuremath{0{,}4232}}
\newcommand{\nmBThreeETTmOneHSevenTwoZeroDLinearMSEMae}{\ensuremath{0{,}4181}}
\newcommand{\nmBThreeETTmOneHSevenTwoZeroDLinearMAEMse}{\ensuremath{0{,}4213}}
\newcommand{\nmBThreeETTmOneHSevenTwoZeroDLinearMAEMae}{\ensuremath{0{,}4101}}
\newcommand{\nmBThreeETTmOneHSevenTwoZeroDLinearHuberMse}{\ensuremath{0{,}4208}}
\newcommand{\nmBThreeETTmOneHSevenTwoZeroDLinearHuberMae}{\ensuremath{0{,}4121}}
\newcommand{\nmBThreeETTmOneHSevenTwoZeroDLinearMSELam}{\ensuremath{2512}}
\newcommand{\nmDeltaETTmOneHSevenTwoZeroDLinear}{\ensuremath{0{,}589}}
\newcommand{\nmBThreeETTmOneHSevenTwoZeroNLinearMSEMse}{\ensuremath{0{,}4262}}
\newcommand{\nmBThreeETTmOneHSevenTwoZeroNLinearMSEMae}{\ensuremath{0{,}4157}}
\newcommand{\nmBThreeETTmOneHSevenTwoZeroNLinearMAEMse}{\ensuremath{0{,}4270}}
\newcommand{\nmBThreeETTmOneHSevenTwoZeroNLinearMAEMae}{\ensuremath{0{,}4113}}
\newcommand{\nmBThreeETTmOneHSevenTwoZeroNLinearHuberMse}{\ensuremath{0{,}4264}}
\newcommand{\nmBThreeETTmOneHSevenTwoZeroNLinearHuberMae}{\ensuremath{0{,}4126}}
\newcommand{\nmBThreeETTmOneHSevenTwoZeroNLinearMSELam}{\ensuremath{7079}}
\newcommand{\nmDeltaETTmOneHSevenTwoZeroNLinear}{\ensuremath{0{,}589}}
\newcommand{\nmBThreeBestMSEMseCount}{5}
\newcommand{\nmBThreeBestMAEMseCount}{24}
\newcommand{\nmBThreeBestHuberMseCount}{7}
\newcommand{\nmBThreeHuberRankOneMseCount}{7}
\newcommand{\nmBThreeHuberRankTwoMseCount}{28}
\newcommand{\nmBThreeHuberRankThreeMseCount}{1}
\newcommand{\nmBThreeBestMSEMaeCount}{1}
\newcommand{\nmBThreeBestMAEMaeCount}{33}
\newcommand{\nmBThreeBestHuberMaeCount}{2}
\newcommand{\nmBThreeHuberRankOneMaeCount}{2}
\newcommand{\nmBThreeHuberRankTwoMaeCount}{34}
\newcommand{\nmBThreeHuberRankThreeMaeCount}{0}
\newcommand{\nmBThreeCellCount}{36}
\newcommand{\nmBThreeMaeBeatsMseETThOneLinearCount}{4}
\newcommand{\nmBThreeMaeBeatsMseETThOneDLinearCount}{4}
\newcommand{\nmBThreeMaeBeatsMseETThOneNLinearCount}{3}
\newcommand{\nmBThreeMaeGainETThOneLinDLinRange}{\ensuremath{0{,}001}--\ensuremath{0{,}015}}
\newcommand{\nmBThreeMaeBeatsMseETThTwoLinearCount}{4}
\newcommand{\nmBThreeMaeBeatsMseETThTwoDLinearCount}{4}
\newcommand{\nmBThreeMaeBeatsMseETThTwoNLinearCount}{0}
\newcommand{\nmBThreeMaeGainETThTwoLinDLinRange}{\ensuremath{0{,}023}--\ensuremath{0{,}151}}
\newcommand{\nmBThreeMaeBeatsMseETTmOneLinearCount}{4}
\newcommand{\nmBThreeMaeBeatsMseETTmOneDLinearCount}{4}
\newcommand{\nmBThreeMaeBeatsMseETTmOneNLinearCount}{3}
\newcommand{\nmBThreeMaeGainETTmOneLinDLinRange}{\ensuremath{0{,}002}--\ensuremath{0{,}006}}
\newcommand{\nmBThreeMaeBeatsMseTotalCount}{30}
\newcommand{\nmBThreeMaeGainETThTwoNLinearRange}{\ensuremath{-0{,}0178}--\ensuremath{-0{,}0002}}
\newcommand{\nmDeltaMin}{\ensuremath{0{,}37}}
\newcommand{\nmDeltaMax}{\ensuremath{0{,}73}}
\newcommand{\nmDeltaMedian}{\ensuremath{0{,}54}}
\newcommand{\nmDeltaRatioMin}{\ensuremath{1{,}4}}
\newcommand{\nmDeltaRatioMax}{\ensuremath{2{,}7}}
\newcommand{\nmDeltaOutlierRange}{\ensuremath{0{,}55}--\ensuremath{1{,}11}}
\newcommand{\nmDeltaMadBetterMseCount}{27}
\newcommand{\nmDeltaMadBetterMaeCount}{34}
\newcommand{\nmDeltaETThTwoLinearHSevenTwoZeroOne}{\ensuremath{0{,}6662}}
\newcommand{\nmDeltaETThTwoLinearHSevenTwoZeroMad}{\ensuremath{0{,}6277}}
\newcommand{\nmDeltaSensETThTwoLinMseRange}{\ensuremath{0{,}602}--\ensuremath{0{,}735}}
\newcommand{\nmDeltaSensRankTwoCount}{50}
\newcommand{\nmDeltaSensCount}{60}
\newcommand{\nmDeltaSensETThOneAmpRange}{\ensuremath{0{,}003}--\ensuremath{0{,}012}}
\newcommand{\nmDeltaSensMaeMonotoneCount}{6}
\newcommand{\nmOutTargetMSEMax}{\ensuremath{0{,}014}}
\newcommand{\nmOutTargetRobustAbsMax}{\ensuremath{0{,}003}}
\newcommand{\nmOutSeriesBestMSEROneCount}{0}
\newcommand{\nmOutSeriesBestMAEROneCount}{8}
\newcommand{\nmOutSeriesBestHuberROneCount}{0}
\newcommand{\nmOutSeriesBestMSERFiveCount}{7}
\newcommand{\nmOutSeriesBestMAERFiveCount}{0}
\newcommand{\nmOutSeriesBestHuberRFiveCount}{1}
\newcommand{\nmOutSeriesBestMSEROneZeroCount}{8}
\newcommand{\nmOutSeriesBestMAEROneZeroCount}{0}
\newcommand{\nmOutSeriesBestHuberROneZeroCount}{0}
\newcommand{\nmOutSeriesDegMSERange}{\ensuremath{0{,}059}--\ensuremath{0{,}077}}
\newcommand{\nmOutSeriesDegMAERange}{\ensuremath{0{,}076}--\ensuremath{0{,}114}}
\newcommand{\nmOutSeriesDegHuberRange}{\ensuremath{0{,}067}--\ensuremath{0{,}103}}
\newcommand{\nmOutSeriesDegOrderCount}{8}
\newcommand{\nmOutSeedSpreadRange}{\ensuremath{0{,}012}--\ensuremath{0{,}015}}
\newcommand{\nmOutLamZeroCount}{36}
\newcommand{\nmOutLamZeroMseMaeCount}{64}
\newcommand{\nmOutLamZeroMseMaeTotalCount}{64}
\newcommand{\nmMagHNineSixSeriesRTenMThreeMSE}{\ensuremath{0{,}3799}}
\newcommand{\nmMagHNineSixSeriesRTenMThreeMAE}{\ensuremath{0{,}3745}}
\newcommand{\nmMagHNineSixSeriesRTenMThreeHuber}{\ensuremath{0{,}3757}}
\newcommand{\nmMagHNineSixSeriesRTenMFiveMSE}{\ensuremath{0{,}3954}}
\newcommand{\nmMagHNineSixSeriesRTenMFiveMAE}{\ensuremath{0{,}3969}}
\newcommand{\nmMagHNineSixSeriesRTenMFiveHuber}{\ensuremath{0{,}3958}}
\newcommand{\nmMagHSevenTwoZeroSeriesRTenMThreeMSE}{\ensuremath{0{,}4757}}
\newcommand{\nmMagHSevenTwoZeroSeriesRTenMThreeMAE}{\ensuremath{0{,}4627}}
\newcommand{\nmMagHSevenTwoZeroSeriesRTenMThreeHuber}{\ensuremath{0{,}4644}}
\newcommand{\nmMagHSevenTwoZeroSeriesRTenMFiveMSE}{\ensuremath{0{,}4904}}
\newcommand{\nmMagHSevenTwoZeroSeriesRTenMFiveMAE}{\ensuremath{0{,}4814}}
\newcommand{\nmMagHSevenTwoZeroSeriesRTenMFiveHuber}{\ensuremath{0{,}4830}}
\newcommand{\nmMagHNineSixSeriesROneMThreeMSEDeg}{\ensuremath{+0{,}0096}}
\newcommand{\nmMagHNineSixSeriesROneMThreeMAEDeg}{\ensuremath{+0{,}0091}}
\newcommand{\nmMagHNineSixSeriesROneMThreeHuberDeg}{\ensuremath{+0{,}0099}}
\newcommand{\nmMagHNineSixFiveSpread}{\ensuremath{0{,}0015}}
\newcommand{\nmMagTargetMSEDegMThree}{\ensuremath{+0{,}0019}}
\newcommand{\nmMagTargetMSEDegMFive}{\ensuremath{+0{,}0047}}
\newcommand{\nmMagTargetMSEDegMOneZero}{\ensuremath{+0{,}0139}}
\newcommand{\nmBOneInsideCount}{28}
\newcommand{\nmBOneOutsideCount}{8}
\newcommand{\nmBOneOutsideDLinearCount}{5}
\newcommand{\nmBOneOutsideLinearCount}{0}
\newcommand{\nmBOneOutsideBelowCount}{8}
\newcommand{\nmBOneETThOneHNineSixLinearMean}{\ensuremath{0{,}3786}}
\newcommand{\nmBOneETThOneHNineSixLinearSd}{\ensuremath{0{,}0107}}
\newcommand{\nmBOneETThOneHNineSixLinearPub}{\ensuremath{0{,}375}}
\newcommand{\nmBOneETThOneHNineSixLinearN}{20}
\newcommand{\nmBOneETThOneHNineSixDLinearMean}{\ensuremath{0{,}3766}}
\newcommand{\nmBOneETThOneHNineSixDLinearSd}{\ensuremath{0{,}0057}}
\newcommand{\nmBOneETThOneHNineSixDLinearPub}{\ensuremath{0{,}375}}
\newcommand{\nmBOneETThOneHNineSixDLinearN}{20}
\newcommand{\nmBOneETThOneHNineSixNLinearMean}{\ensuremath{0{,}3779}}
\newcommand{\nmBOneETThOneHNineSixNLinearSd}{\ensuremath{0{,}0072}}
\newcommand{\nmBOneETThOneHNineSixNLinearPub}{\ensuremath{0{,}374}}
\newcommand{\nmBOneETThOneHNineSixNLinearN}{20}
\newcommand{\nmBOneETThOneHOneNineTwoLinearMean}{\ensuremath{0{,}4250}}
\newcommand{\nmBOneETThOneHOneNineTwoLinearSd}{\ensuremath{0{,}0222}}
\newcommand{\nmBOneETThOneHOneNineTwoLinearPub}{\ensuremath{0{,}418}}
\newcommand{\nmBOneETThOneHOneNineTwoLinearN}{20}
\newcommand{\nmBOneETThOneHOneNineTwoDLinearMean}{\ensuremath{0{,}4146}}
\newcommand{\nmBOneETThOneHOneNineTwoDLinearSd}{\ensuremath{0{,}0152}}
\newcommand{\nmBOneETThOneHOneNineTwoDLinearPub}{\ensuremath{0{,}405}}
\newcommand{\nmBOneETThOneHOneNineTwoDLinearN}{20}
\newcommand{\nmBOneETThOneHOneNineTwoNLinearMean}{\ensuremath{0{,}4105}}
\newcommand{\nmBOneETThOneHOneNineTwoNLinearSd}{\ensuremath{0{,}0106}}
\newcommand{\nmBOneETThOneHOneNineTwoNLinearPub}{\ensuremath{0{,}408}}
\newcommand{\nmBOneETThOneHOneNineTwoNLinearN}{20}
\newcommand{\nmBOneETThOneHThreeThreeSixLinearMean}{\ensuremath{0{,}4553}}
\newcommand{\nmBOneETThOneHThreeThreeSixLinearSd}{\ensuremath{0{,}0189}}
\newcommand{\nmBOneETThOneHThreeThreeSixLinearPub}{\ensuremath{0{,}479}}
\newcommand{\nmBOneETThOneHThreeThreeSixLinearN}{20}
\newcommand{\nmBOneETThOneHThreeThreeSixDLinearMean}{\ensuremath{0{,}4573}}
\newcommand{\nmBOneETThOneHThreeThreeSixDLinearSd}{\ensuremath{0{,}0322}}
\newcommand{\nmBOneETThOneHThreeThreeSixDLinearPub}{\ensuremath{0{,}439}}
\newcommand{\nmBOneETThOneHThreeThreeSixDLinearN}{20}
\newcommand{\nmBOneETThOneHThreeThreeSixNLinearMean}{\ensuremath{0{,}4320}}
\newcommand{\nmBOneETThOneHThreeThreeSixNLinearSd}{\ensuremath{0{,}0042}}
\newcommand{\nmBOneETThOneHThreeThreeSixNLinearPub}{\ensuremath{0{,}429}}
\newcommand{\nmBOneETThOneHThreeThreeSixNLinearN}{20}
\newcommand{\nmBOneETThOneHSevenTwoZeroLinearMean}{\ensuremath{0{,}5063}}
\newcommand{\nmBOneETThOneHSevenTwoZeroLinearSd}{\ensuremath{0{,}0363}}
\newcommand{\nmBOneETThOneHSevenTwoZeroLinearPub}{\ensuremath{0{,}624}}
\newcommand{\nmBOneETThOneHSevenTwoZeroLinearN}{20}
\newcommand{\nmBOneETThOneHSevenTwoZeroDLinearMean}{\ensuremath{0{,}5096}}
\newcommand{\nmBOneETThOneHSevenTwoZeroDLinearSd}{\ensuremath{0{,}0396}}
\newcommand{\nmBOneETThOneHSevenTwoZeroDLinearPub}{\ensuremath{0{,}472}}
\newcommand{\nmBOneETThOneHSevenTwoZeroDLinearN}{20}
\newcommand{\nmBOneETThOneHSevenTwoZeroNLinearMean}{\ensuremath{0{,}4363}}
\newcommand{\nmBOneETThOneHSevenTwoZeroNLinearSd}{\ensuremath{0{,}0014}}
\newcommand{\nmBOneETThOneHSevenTwoZeroNLinearPub}{\ensuremath{0{,}440}}
\newcommand{\nmBOneETThOneHSevenTwoZeroNLinearN}{20}
\newcommand{\nmBOneETThTwoHNineSixLinearMean}{\ensuremath{0{,}2968}}
\newcommand{\nmBOneETThTwoHNineSixLinearSd}{\ensuremath{0{,}0191}}
\newcommand{\nmBOneETThTwoHNineSixLinearPub}{\ensuremath{0{,}288}}
\newcommand{\nmBOneETThTwoHNineSixLinearN}{20}
\newcommand{\nmBOneETThTwoHNineSixDLinearMean}{\ensuremath{0{,}2913}}
\newcommand{\nmBOneETThTwoHNineSixDLinearSd}{\ensuremath{0{,}0060}}
\newcommand{\nmBOneETThTwoHNineSixDLinearPub}{\ensuremath{0{,}289}}
\newcommand{\nmBOneETThTwoHNineSixDLinearN}{20}
\newcommand{\nmBOneETThTwoHNineSixNLinearMean}{\ensuremath{0{,}2770}}
\newcommand{\nmBOneETThTwoHNineSixNLinearSd}{\ensuremath{0{,}0020}}
\newcommand{\nmBOneETThTwoHNineSixNLinearPub}{\ensuremath{0{,}277}}
\newcommand{\nmBOneETThTwoHNineSixNLinearN}{10}
\newcommand{\nmBOneETThTwoHOneNineTwoLinearMean}{\ensuremath{0{,}3961}}
\newcommand{\nmBOneETThTwoHOneNineTwoLinearSd}{\ensuremath{0{,}0472}}
\newcommand{\nmBOneETThTwoHOneNineTwoLinearPub}{\ensuremath{0{,}377}}
\newcommand{\nmBOneETThTwoHOneNineTwoLinearN}{20}
\newcommand{\nmBOneETThTwoHOneNineTwoDLinearMean}{\ensuremath{0{,}3807}}
\newcommand{\nmBOneETThTwoHOneNineTwoDLinearSd}{\ensuremath{0{,}0145}}
\newcommand{\nmBOneETThTwoHOneNineTwoDLinearPub}{\ensuremath{0{,}383}}
\newcommand{\nmBOneETThTwoHOneNineTwoDLinearN}{20}
\newcommand{\nmBOneETThTwoHOneNineTwoNLinearMean}{\ensuremath{0{,}3444}}
\newcommand{\nmBOneETThTwoHOneNineTwoNLinearSd}{\ensuremath{0{,}0041}}
\newcommand{\nmBOneETThTwoHOneNineTwoNLinearPub}{\ensuremath{0{,}344}}
\newcommand{\nmBOneETThTwoHOneNineTwoNLinearN}{10}
\newcommand{\nmBOneETThTwoHThreeThreeSixLinearMean}{\ensuremath{0{,}4496}}
\newcommand{\nmBOneETThTwoHThreeThreeSixLinearSd}{\ensuremath{0{,}0241}}
\newcommand{\nmBOneETThTwoHThreeThreeSixLinearPub}{\ensuremath{0{,}452}}
\newcommand{\nmBOneETThTwoHThreeThreeSixLinearN}{20}
\newcommand{\nmBOneETThTwoHThreeThreeSixDLinearMean}{\ensuremath{0{,}4529}}
\newcommand{\nmBOneETThTwoHThreeThreeSixDLinearSd}{\ensuremath{0{,}0242}}
\newcommand{\nmBOneETThTwoHThreeThreeSixDLinearPub}{\ensuremath{0{,}448}}
\newcommand{\nmBOneETThTwoHThreeThreeSixDLinearN}{20}
\newcommand{\nmBOneETThTwoHThreeThreeSixNLinearMean}{\ensuremath{0{,}3804}}
\newcommand{\nmBOneETThTwoHThreeThreeSixNLinearSd}{\ensuremath{0{,}0139}}
\newcommand{\nmBOneETThTwoHThreeThreeSixNLinearPub}{\ensuremath{0{,}357}}
\newcommand{\nmBOneETThTwoHThreeThreeSixNLinearN}{10}
\newcommand{\nmBOneETThTwoHSevenTwoZeroLinearMean}{\ensuremath{0{,}6967}}
\newcommand{\nmBOneETThTwoHSevenTwoZeroLinearSd}{\ensuremath{0{,}0484}}
\newcommand{\nmBOneETThTwoHSevenTwoZeroLinearPub}{\ensuremath{0{,}698}}
\newcommand{\nmBOneETThTwoHSevenTwoZeroLinearN}{20}
\newcommand{\nmBOneETThTwoHSevenTwoZeroDLinearMean}{\ensuremath{0{,}6869}}
\newcommand{\nmBOneETThTwoHSevenTwoZeroDLinearSd}{\ensuremath{0{,}0592}}
\newcommand{\nmBOneETThTwoHSevenTwoZeroDLinearPub}{\ensuremath{0{,}605}}
\newcommand{\nmBOneETThTwoHSevenTwoZeroDLinearN}{20}
\newcommand{\nmBOneETThTwoHSevenTwoZeroNLinearMean}{\ensuremath{0{,}4058}}
\newcommand{\nmBOneETThTwoHSevenTwoZeroNLinearSd}{\ensuremath{0{,}0122}}
\newcommand{\nmBOneETThTwoHSevenTwoZeroNLinearPub}{\ensuremath{0{,}394}}
\newcommand{\nmBOneETThTwoHSevenTwoZeroNLinearN}{10}
\newcommand{\nmBOneETTmOneHNineSixLinearMean}{\ensuremath{0{,}3057}}
\newcommand{\nmBOneETTmOneHNineSixLinearSd}{\ensuremath{0{,}0019}}
\newcommand{\nmBOneETTmOneHNineSixLinearPub}{\ensuremath{0{,}308}}
\newcommand{\nmBOneETTmOneHNineSixLinearN}{20}
\newcommand{\nmBOneETTmOneHNineSixDLinearMean}{\ensuremath{0{,}3017}}
\newcommand{\nmBOneETTmOneHNineSixDLinearSd}{\ensuremath{0{,}0018}}
\newcommand{\nmBOneETTmOneHNineSixDLinearPub}{\ensuremath{0{,}299}}
\newcommand{\nmBOneETTmOneHNineSixDLinearN}{20}
\newcommand{\nmBOneETTmOneHNineSixNLinearMean}{\ensuremath{0{,}3091}}
\newcommand{\nmBOneETTmOneHNineSixNLinearSd}{\ensuremath{0{,}0042}}
\newcommand{\nmBOneETTmOneHNineSixNLinearPub}{\ensuremath{0{,}306}}
\newcommand{\nmBOneETTmOneHNineSixNLinearN}{20}
\newcommand{\nmBOneETTmOneHOneNineTwoLinearMean}{\ensuremath{0{,}3401}}
\newcommand{\nmBOneETTmOneHOneNineTwoLinearSd}{\ensuremath{0{,}0012}}
\newcommand{\nmBOneETTmOneHOneNineTwoLinearPub}{\ensuremath{0{,}340}}
\newcommand{\nmBOneETTmOneHOneNineTwoLinearN}{20}
\newcommand{\nmBOneETTmOneHOneNineTwoDLinearMean}{\ensuremath{0{,}3369}}
\newcommand{\nmBOneETTmOneHOneNineTwoDLinearSd}{\ensuremath{0{,}0015}}
\newcommand{\nmBOneETTmOneHOneNineTwoDLinearPub}{\ensuremath{0{,}335}}
\newcommand{\nmBOneETTmOneHOneNineTwoDLinearN}{20}
\newcommand{\nmBOneETTmOneHOneNineTwoNLinearMean}{\ensuremath{0{,}3436}}
\newcommand{\nmBOneETTmOneHOneNineTwoNLinearSd}{\ensuremath{0{,}0028}}
\newcommand{\nmBOneETTmOneHOneNineTwoNLinearPub}{\ensuremath{0{,}349}}
\newcommand{\nmBOneETTmOneHOneNineTwoNLinearN}{20}
\newcommand{\nmBOneETTmOneHThreeThreeSixLinearMean}{\ensuremath{0{,}3769}}
\newcommand{\nmBOneETTmOneHThreeThreeSixLinearSd}{\ensuremath{0{,}0025}}
\newcommand{\nmBOneETTmOneHThreeThreeSixLinearPub}{\ensuremath{0{,}376}}
\newcommand{\nmBOneETTmOneHThreeThreeSixLinearN}{20}
\newcommand{\nmBOneETTmOneHThreeThreeSixDLinearMean}{\ensuremath{0{,}3731}}
\newcommand{\nmBOneETTmOneHThreeThreeSixDLinearSd}{\ensuremath{0{,}0027}}
\newcommand{\nmBOneETTmOneHThreeThreeSixDLinearPub}{\ensuremath{0{,}369}}
\newcommand{\nmBOneETTmOneHThreeThreeSixDLinearN}{20}
\newcommand{\nmBOneETTmOneHThreeThreeSixNLinearMean}{\ensuremath{0{,}3779}}
\newcommand{\nmBOneETTmOneHThreeThreeSixNLinearSd}{\ensuremath{0{,}0022}}
\newcommand{\nmBOneETTmOneHThreeThreeSixNLinearPub}{\ensuremath{0{,}375}}
\newcommand{\nmBOneETTmOneHThreeThreeSixNLinearN}{20}
\newcommand{\nmBOneETTmOneHSevenTwoZeroLinearMean}{\ensuremath{0{,}4336}}
\newcommand{\nmBOneETTmOneHSevenTwoZeroLinearSd}{\ensuremath{0{,}0039}}
\newcommand{\nmBOneETTmOneHSevenTwoZeroLinearPub}{\ensuremath{0{,}440}}
\newcommand{\nmBOneETTmOneHSevenTwoZeroLinearN}{20}
\newcommand{\nmBOneETTmOneHSevenTwoZeroDLinearMean}{\ensuremath{0{,}4289}}
\newcommand{\nmBOneETTmOneHSevenTwoZeroDLinearSd}{\ensuremath{0{,}0028}}
\newcommand{\nmBOneETTmOneHSevenTwoZeroDLinearPub}{\ensuremath{0{,}425}}
\newcommand{\nmBOneETTmOneHSevenTwoZeroDLinearN}{20}
\newcommand{\nmBOneETTmOneHSevenTwoZeroNLinearMean}{\ensuremath{0{,}4346}}
\newcommand{\nmBOneETTmOneHSevenTwoZeroNLinearSd}{\ensuremath{0{,}0031}}
\newcommand{\nmBOneETTmOneHSevenTwoZeroNLinearPub}{\ensuremath{0{,}433}}
\newcommand{\nmBOneETTmOneHSevenTwoZeroNLinearN}{20}
\newcommand{\nmBFourSgdBetterCount}{6}
\newcommand{\nmBFourCfBetterCount}{10}
\newcommand{\nmBFourValTestCount}{7}
\newcommand{\nmBFourSgdOverfitCount}{12}
\newcommand{\nmBFourCfOverfitCount}{0}
\newcommand{\nmBFourTieCount}{1}
\newcommand{\nmBFourETThOneDTestRange}{\ensuremath{+0{,}0034}--\ensuremath{+0{,}0400}}
\newcommand{\nmBFourETThOneWorseCount}{12}
\newcommand{\nmBFourETTmOneDTestRange}{\ensuremath{+0{,}0025}--\ensuremath{+0{,}0104}}
\newcommand{\nmBFourETTmOneWorseCount}{12}
\newcommand{\nmTunedSgdBetterCount}{5}
\newcommand{\nmTunedSgdBetterXOneCount}{6}
\newcommand{\nmTunedCfBetterCount}{2}
\newcommand{\nmTunedCfBetterXOneCount}{10}
\newcommand{\nmTunedValTestCount}{8}
\newcommand{\nmTunedValTestXOneCount}{7}
\newcommand{\nmTunedSgdOverfitCount}{17}
\newcommand{\nmTunedSgdOverfitXOneCount}{11}
\newcommand{\nmTunedCfOverfitCount}{0}
\newcommand{\nmTunedCfOverfitXOneCount}{0}
\newcommand{\nmTunedTieCount}{4}
\newcommand{\nmTunedTieXOneCount}{2}
\newcommand{\nmTunedWorseSigCount}{27}
\newcommand{\nmTunedWorseMeanCount}{28}
\newcommand{\nmTunedChangedCount}{16}
\newcommand{\nmTunedChangedByLrCount}{14}
\newcommand{\nmTunedChosenOneCount}{7}
\newcommand{\nmTunedChosenFourCount}{7}
\newcommand{\nmTunedChosenZeroPointFiveCount}{6}
\newcommand{\nmTunedChosenZeroPointTwoFiveCount}{1}
\newcommand{\nmTunedChosenZeroPointOneTwoFiveCount}{1}
\newcommand{\nmTunedChosenEightCount}{5}
\newcommand{\nmTunedChosenTwoCount}{7}
\newcommand{\nmTunedChosenThreeTwoCount}{1}
\newcommand{\nmTunedChosenOneSixCount}{1}
\newcommand{\nmTunedHMOneWorseSigCount}{23}
\newcommand{\nmTunedHMOneWorseMeanCount}{24}
\newcommand{\nmTunedHMOneCellCount}{24}
\newcommand{\nmTunedHMOneDTestRange}{\ensuremath{+0{,}0020}--\ensuremath{+0{,}0290}}
\newcommand{\nmTunedLrEdgeOldCount}{15}
\newcommand{\nmTunedLrEdgeCount}{2}
\newcommand{\nmTunedLrExtChangedCount}{4}
\newcommand{\nmTunedLrExtWallHours}{\ensuremath{0{,}8}}
\newcommand{\nmTunedLrChosenList}{$\times 0{,}125$: 1 ô; $\times 0{,}25$: 1; $\times 0{,}5$: 6; $\times 1$: 7; $\times 2$: 7; $\times 4$: 7; $\times 8$: 5; $\times 16$: 1; $\times 32$: 1}
\newcommand{\nmTunedLrStillEdgeList}{ETTm1 Linear $H = 96$ ở $\times 32$; ETTh2 NLinear $H = 720$ ở $\times 0{,}125$}
\newcommand{\nmTunedLrNoisyCount}{31}
\newcommand{\nmTunedLrNoisyMedianRatio}{\ensuremath{0{,}37}}
\newcommand{\nmTunedLrCellCount}{36}
\newcommand{\nmTunedETTmOneLinearChoseGeEightCount}{3}
\newcommand{\nmTunedETTmOneLinearChoseLeTwoCount}{1}
\newcommand{\nmTunedETTmOneDLinearChoseGeEightCount}{0}
\newcommand{\nmTunedETTmOneDLinearChoseLeTwoCount}{3}
\newcommand{\nmCThreeDLinBetterBothCount}{4}
\newcommand{\nmCThreeLinBetterTestCount}{0}
\newcommand{\nmCThreeSigDTestRange}{\ensuremath{-0{,}0032}--\ensuremath{-0{,}0047}}
\newcommand{\nmCThreeETThOneHSevenTwoZeroDTest}{\ensuremath{+0{,}003}}
\newcommand{\nmCThreeETThOneHSevenTwoZeroDTestLo}{\ensuremath{-0{,}021}}
\newcommand{\nmCThreeETThOneHSevenTwoZeroDTestHi}{\ensuremath{0{,}028}}
\newcommand{\nmCThreeETThOneHSevenTwoZeroPubD}{\ensuremath{-0{,}152}}
\newcommand{\nmCThreeETThTwoHSevenTwoZeroDTest}{\ensuremath{-0{,}010}}
\newcommand{\nmCThreeETThTwoHSevenTwoZeroDTestLo}{\ensuremath{-0{,}044}}
\newcommand{\nmCThreeETThTwoHSevenTwoZeroDTestHi}{\ensuremath{0{,}025}}
\newcommand{\nmCThreeETThTwoHSevenTwoZeroPubD}{\ensuremath{-0{,}093}}
\newcommand{\nmTunedCThreeDLinBetterCount}{0}
\newcommand{\nmTunedCThreeXOneETTmOneTestCount}{4}
\newcommand{\nmTunedCThreeETTmOneDTestRange}{\ensuremath{-0{,}0020}--\ensuremath{0{,}0004}}
\newcommand{\nmTunedCThreeETTmOneCIHalfMax}{\ensuremath{0{,}0073}}
\newcommand{\nmEThTwoCfFullHSevenTwoZero}{\ensuremath{0{,}740}}
\newcommand{\nmEThTwoCfFiltHSevenTwoZero}{\ensuremath{0{,}646}}
\newcommand{\nmEThTwoSgdFullHSevenTwoZero}{\ensuremath{0{,}697}}
\newcommand{\nmEThTwoSgdFiltHSevenTwoZero}{\ensuremath{0{,}601}}
\newcommand{\nmEThTwoFiltGainHSevenTwoZeroRange}{\ensuremath{0{,}076}--\ensuremath{0{,}095}}
\newcommand{\nmEThTwoFiltLossHNineSixRange}{\ensuremath{0{,}01}--\ensuremath{0{,}02}}
\newcommand{\nmEThTwoMaeGainRemainPctHThreeThreeSix}{47}
\newcommand{\nmEThTwoMaeGainRemainPctHSevenTwoZero}{60}
\newcommand{\nmEThTwoConstTrainPctRange}{\ensuremath{2{,}5}--\ensuremath{6{,}3}}
\newcommand{\nmEThTwoSgdStillBetterCount}{3}
\newcommand{\nmDropPubInTrueCount}{1}
\newcommand{\nmDropPubInFalseCount}{0}
\newcommand{\nmDropCellCount}{8}
\newcommand{\nmDropDiffAbsMax}{\ensuremath{5\cdot 10^{-3}}}
\newcommand{\nmDropDroppedRange}{1--17}
\newcommand{\nmDropDroppedPctMax}{\ensuremath{0{,}79}}
\newcommand{\nmDropPubBelowTrueCount}{7}
\newcommand{\nmDropPubGapRange}{\ensuremath{0{,}0004}--\ensuremath{0{,}0189}}
\newcommand{\nmDropSeedCount}{5}
\newcommand{\nmOffETTmOneHNineSixMean}{\ensuremath{0{,}3018}}
\newcommand{\nmOffETTmOneHNineSixSd}{\ensuremath{0{,}0010}}
\newcommand{\nmOffETTmOneHNineSixN}{5}
\newcommand{\nmRidgeBestWorseOrigCount}{8}
\newcommand{\nmRidgeBestWorseTunedCount}{8}
\newcommand{\nmRidgeCellCount}{8}
\newcommand{\nmRidgeGapOrigRange}{\ensuremath{0{,}0005}--\ensuremath{0{,}0534}}
\newcommand{\nmRidgeGapTunedRange}{\ensuremath{0{,}0020}--\ensuremath{0{,}0507}}
\newcommand{\nmRidgeStarMinusBestMax}{\ensuremath{0{,}0001}}
\newcommand{\nmStuckRunCount}{163}
\newcommand{\nmStuckRunLongCount}{19}
\newcommand{\nmStuckDayThirtyOneCount}{63}
\newcommand{\nmStuckRunMaxHours}{1025}
\newcommand{\nmStuckRunLongTrainCount}{10}
\newcommand{\nmStuckRunLongValCount}{1}
\newcommand{\nmStuckRunLongTestCount}{8}
\newcommand{\nmDeltaHRankChangedCount}{0}
\newcommand{\nmDeltaHCellCount}{2}
\newcommand{\nmDeltaHETThOneHNineSixLinearMin}{\ensuremath{0{,}243}}
\newcommand{\nmDeltaHETThOneHNineSixLinearMax}{\ensuremath{0{,}595}}
\newcommand{\nmDeltaHETThOneHNineSixLinearMad}{\ensuremath{0{,}538}}
\newcommand{\nmDeltaHETThOneHNineSixLinearHubMadMse}{\ensuremath{0{,}3658}}
\newcommand{\nmDeltaHETThOneHNineSixLinearHubMadMae}{\ensuremath{0{,}3852}}
\newcommand{\nmDeltaHETThOneHNineSixLinearHubHMse}{\ensuremath{0{,}3656}}
\newcommand{\nmDeltaHETThOneHNineSixLinearHubHMae}{\ensuremath{0{,}3851}}
\newcommand{\nmDeltaHETThTwoHSevenTwoZeroLinearMin}{\ensuremath{0{,}174}}
\newcommand{\nmDeltaHETThTwoHSevenTwoZeroLinearMax}{\ensuremath{0{,}801}}
\newcommand{\nmDeltaHETThTwoHSevenTwoZeroLinearMad}{\ensuremath{0{,}610}}
\newcommand{\nmDeltaHETThTwoHSevenTwoZeroLinearHubMadMse}{\ensuremath{0{,}6277}}
\newcommand{\nmDeltaHETThTwoHSevenTwoZeroLinearHubMadMae}{\ensuremath{0{,}5541}}
\newcommand{\nmDeltaHETThTwoHSevenTwoZeroLinearHubHMse}{\ensuremath{0{,}6357}}
\newcommand{\nmDeltaHETThTwoHSevenTwoZeroLinearHubHMae}{\ensuremath{0{,}5566}}
\newcommand{\nmReproDeltaHMinutes}{3}
\newcommand{\nmReproDropGpuHours}{\ensuremath{2{,}6}}

% Mục từ của bảng thuật ngữ (glossaries noidx: phải định nghĩa trong preamble). In ở sections/00_thuat_ngu.tex.
\newglossaryentry{closedform}{name={closed-form solution}, user1={nghiệm dạng đóng}, description={Nghiệm tính trực tiếp bằng công thức đại số (OLS, ridge, NLinear có ràng buộc), không lặp.}}
\newglossaryentry{optimum}{name={optimum}, user1={nghiệm tối ưu}, description={Điểm cực tiểu toàn cục của hàm mục tiêu trên tập huấn luyện. Với MSE là nghiệm dạng đóng; với MAE, Huber là điểm IRLS hội tụ tới.}}
\newglossaryentry{objective}{name={objective}, user1={hàm mục tiêu}, description={Hàm được cực tiểu khi huấn luyện: MSE, MAE hoặc Huber, cộng phạt $\lambda\|W\|^2$ nếu có.}}
\newglossaryentry{val}{name={validation set}, user1={tập kiểm định (val)}, description={Dùng để chọn $\lambda$, $\delta$, learning rate và epoch. Không bao giờ dùng tập kiểm tra để chọn.}}
\newglossaryentry{test}{name={test set}, user1={tập kiểm tra (test)}, description={Chỉ dùng để báo cáo sai số cuối.}}
\newglossaryentry{stuck}{name={stuck sensor segment}, user1={đoạn cảm biến bị kẹt}, description={Đoạn chuỗi giữ nguyên một giá trị trong thời gian dài (ở ETTh2); cửa sổ có đầu vào hoặc đích hằng.}}
\newglossaryentry{horizon}{name={horizon $H$}, user1={độ dài dự báo}, description={Số bước tương lai cần dự báo: 96, 192, 336, 720.}}
\newglossaryentry{lookback}{name={look-back window $L$}, user1={cửa sổ đầu vào}, description={Số bước quá khứ làm đầu vào; cố định $L = 336$.}}
\newglossaryentry{ridge}{name={ridge}, user1={hồi quy phạt L2}, description={Bình phương tối thiểu cộng $\lambda\|W\|^2$; tương đương weight decay.}}
\newglossaryentry{ols}{name={OLS}, user1={bình phương tối thiểu thường}, description={Ridge với $\lambda = 0$.}}
\newglossaryentry{irls}{name={IRLS}, user1={bình phương tối thiểu có trọng số lặp lại}, description={Thuật toán cực tiểu MAE, Huber bằng một chuỗi bài toán bình phương tối thiểu có trọng số; là thuật toán MM nên hàm mục tiêu không tăng.}}
\newglossaryentry{sgd}{name={SGD}, user1={gradient ngẫu nhiên}, description={Ở đây là cách huấn luyện của LTSF-Linear: Adam theo mini-batch, dừng sớm theo MSE val.}}
\newglossaryentry{seed}{name={seed}, user1={hạt giống ngẫu nhiên}, description={Quyết định khởi tạo và thứ tự mini-batch của SGD, và vị trí outlier nhân tạo.}}
\newglossaryentry{lr}{name={learning rate}, user1={tốc độ học (lr)}, description={Bước của Adam; được tinh chỉnh ở mục~\ref{sec:rq3}.}}
\newglossaryentry{huber}{name={Huber loss}, user1={hàm mất mát Huber}, description={Bậc hai khi $|r| \le \delta$, tuyến tính khi $|r| > \delta$; nằm giữa MSE và MAE.}}
\newglossaryentry{outlier}{name={outlier}, user1={giá trị ngoại lai}, description={Ở đây là điểm được cộng $\pm k$ lần độ lệch chuẩn của kênh trong train.}}

